\subsection{A family with trivial monodromy}\label{sec:centralatoms}
For the centrally symmetric del Pezzo varieties of Section~\ref{sec:example} the entire primitive cover can be split. In dimension four, for instance, the two orbits of cardinality two have isomorphic stabilizers and could in principle be exchanged by monodromy. To distinguish their sheets we use the grading of the quantum product.

\begin{proposition}\label{prop:centralatoms}
For the centrally symmetric del Pezzo variety of dimension $2d$, $d\ge1$, let the full fan group act. The primitive cover is constant as a cover with group action on the entire nonarchimedean fixed base. There are $(d+1)(d+2)/2$ local equivariant Hodge atoms, each of covering degree one. Write $m=2d+1$. For each $d+1\le k\le m$ there is one atom with eigenspace dimension $\binom{m}{k}$ and there are $(2k-m-1)/2$ atoms with eigenspace dimension $2\binom{m}{k}$, counted with these labels if dimensions coincide.
\end{proposition}
\begin{proof}
The character lattice of the divisor torus consists of the integral relations among the rays. Write a relation on the opposite pairs as $(b_i+c,b_i)$, since the only relation among the $v_i$ is their sum. Permuting the pairs gives coinvariant coordinates $(B,c)$ for $B=\sum_i b_i$. Inversion sends these to $(B+mc,-c)$. Thus the full coinvariant lattice is $\Z^2/\Z(m,-2)$. Since $m$ is odd, the anticanonical degree $2B+mc$ identifies this lattice with $\Z$. The entire fixed divisor torus has a unique coordinate $x$ with $q^\beta=x^{c_1\cdot\beta}$. Its valuation domain is $0<|x|<1$. The fixed base, including the invariant higher inputs and the unit line, is connected.

Let $e_i$ denote a homogeneous basis vector of degree $2i$, and use the global frame $e'_i=x^{-i}e_i$, suppressing multiplicity indices. Write $\eta_p$ for the parameters of degree $2p$ with $p\ge2$. A contribution of degree $\beta$ to the $e_k$ coefficient of $e_i\star e_j$, with further inputs of half-degrees $p_1,\ldots,p_l$, satisfies the dimension equation
\[
             c_1\cdot\beta=i+j-k+\sum_{a=1}^l(p_a-1).
\]
Hence in the primed frame the structure constants are power series in $u_p=x^{p-1}\eta_p$ alone. Their coefficients are rational. For prescribed input multiplicities the anticanonical degree is fixed, admitting only finitely many effective curve classes. These formal series define analytic functions on the open unit polydisc since the constant field is trivially valued. All allowed fixed parameters have $|u_p|<1$. The unit coordinate drops from the product. Subtracting its scalar contribution from the Euler element and dividing by $x$ likewise gives an element over this power series ring, with constant term $c_1$ in the normalized small algebra.

At $u=0$ the structure constants in this frame are those of the small product with each Novikov monomial set to one. Its toric presentation sends the divisor of each ray to the corresponding Laurent monomial, hence their sum $c_1$ to $\sum_i(y_i+y_i^{-1})$ \cite[Proposition~3.3]{OT}, \cite[\S1.7]{MM}. This identification respects the fan permutations. The monotone small algebra for this entire family was proved semisimple in \cite[Proposition~1.6.1]{MM}. The equations in Section~\ref{sec:example} give its primitive characters over $\overline{\mathbb Q}$. Each special idempotent lifts uniquely in the finite free algebra over the complete power series ring on the invariant $u$-coordinates. Indeed the linearization of $e^2=e$ at its reduction is multiplication by $2e-1$, an involution. The lifts are orthogonal and sum to one by uniqueness; they form a basis by reduction. Their coefficients, being constants in $\overline{\mathbb Q}$ at each power of $u$, converge on the same open polydisc. Pulling back to the full fixed base gives a simultaneous splitting with the original permutation action.

We check that distinct orbits have different normalized Euler values throughout this base. At $u=0$ the orbit with a majority of $k$ coordinates equal to $r$ has value $j(r+r^{-1})$, where $j=2k-m$ is positive and odd and $r^j=(-1)^{m-k}$. For fixed $j$, equality of two values means that their roots form the same inverse pair. If $r$ is nonreal, $r^2$ has odd order $a>1$ and $r+r^{-1}=r^{-1}(1+r^2)$ is an algebraic unit. In fact squaring permutes the primitive $a$th roots, so cancelling the products of $1-z$ against the products of $1-z^2$ gives $\prod_z(1+z)=1$. For real $r$ the factor is $\pm2$. Equality for two nonreal roots with different $j$ would make the positive rational $j/j'$ an algebraic unit. Two real roots of equal value also force $j=j'$. In the mixed case the same unit argument would make one odd integer twice another. Thus the constant terms separate all distinct orbits. Differences with nonzero constant term remain nonzero at every allowed $u$ by the ultrametric inequality.

There is one real root for each $j$ and there are $(j-1)/2$ nonreal inverse pairs. The group permutes the majority set transitively and simultaneous inversion exchanges the two roots of each nonreal pair. This gives the stated eigenspace dimensions. Summing $(j+1)/2$ over the positive odd numbers $j\le m$ gives the number of separate sheets.
\end{proof}

For the degree-six surface this gives three atoms with eigenspace dimensions $1,2,3$. For the fourfold there are six, with dimensions $1,2,2,5,10,10$. Even the two isomorphic two-element orbits lie on separate components. The ordinary exceptional collection in this family and its fan action are known \cite[Proposition~4.12]{BDM}; see also Castravet--Tevelev \cite[Theorem~6.6]{CTI}. The relation between the three surface blocks and the small spectrum is discussed in \cite[Remark~1.6]{ESS}. Proposition~\ref{prop:centralatoms} computes the degrees on the full fixed quantum base.
